\textbf{结论名称：} $y=A\sin(\omega x+\varphi)$ 的图像变换（平移与伸缩的顺序）

\textbf{适用条件：} $A\neq0$，$\omega>0$，$\varphi\in\mathbb R$，从 $y=\sin x$ 出发。

\textbf{变量说明：} $A$ 是带符号的纵向系数，振幅为 $|A|$；$\omega$ 是正角频率，周期 $T=2\pi/\omega$，横坐标伸缩倍数为 $1/\omega$；$\varphi$ 是初相参数。$\varphi>0$ 时左移，$\varphi<0$ 时右移，$\varphi=0$ 时不移动。

\textbf{核心公式：}
\[
\boxed{A\sin(\omega x+\varphi)=A\sin\!\Big[\omega\Big(x+\frac{\varphi}{\omega}\Big)\Big]}
\]
由 $y=\sin x$ 出发有两条路径：
\begin{enumerate}
\item 先向左平移有符号距离 $\varphi$，再将各点横坐标乘以 $1/\omega$，最后将纵坐标乘以 $A$：$\sin x\to\sin(x+\varphi)\to\sin(\omega x+\varphi)\to A\sin(\omega x+\varphi)$。
\item 先将各点横坐标乘以 $1/\omega$，再向左平移有符号距离 $\varphi/\omega$，最后将纵坐标乘以 $A$：$\sin x\to\sin(\omega x)\to\sin[\omega(x+\varphi/\omega)]\to A\sin(\omega x+\varphi)$。
\end{enumerate}
两条路径的平移距离分别为 $|\varphi|$ 和 $|\varphi|/\omega$；当且仅当 $\varphi=0$ 或 $\omega=1$ 时相等。调整平移量后，两种顺序得到同一图像；一般不可直接交换平移与伸缩。若 $A<0$，最后一步可分为按 $|A|$ 倍纵向伸缩，再关于 $x$ 轴对称。$\omega>1$ 时横向压缩，$\omega=1$ 时不变，$0<\omega<1$ 时横向拉伸。

\textbf{简单例子：} 取 $A=-2$、$\omega=2$、$\varphi=\pi/4$。先移后缩时左移 $\pi/4$ 再横坐标乘以 $1/2$；先缩后移时先横坐标乘以 $1/2$ 再左移 $\pi/8$。两条路径最后都把纵坐标乘以 $-2$，得到 $y=-2\sin(2x+\pi/4)$；在 $x=\pi/8$ 处取值 $-2$。
